保持轨迹形状与$T_f=4\units s$不变，取$\gamma=0^\circ,30^\circ,60^\circ,90^\circ$，计算两电机全程最大绝对力矩。该扫参采用401点，结果依次为0.152、4.180、7.158和$8.247\units{N\,m}$，如图\ref{fig:scenario}(a)所示。由此可见，将同样的描绘机构从水平面转为竖直面，会显著增加本模型的面内驱动负担。这里比较的是同一课程设计机构的安装方向，不是把两种不同Quanser产品的规格作横向对比。

\figfull{fig11_installation_speed.pdf}{描绘场景的安装方向与时间敏感性。(a)同一主轨迹的板面倾角扫描；(b)两条几何轨迹的执行时间扫描，虚线为对应轨迹的静态重力力矩极限。各轨迹使用401点，与主计算1601点的离散峰值可能有微小差异。}{fig:scenario}

对竖直安装，进一步令$T_f=2,4,8,16\units s$，结果见表\ref{tab:speed}。由于$\dot{\vect q}\propto1/T_f$、$\ddot{\vect q}\propto1/T_f^2$，惯性及速度二次项随$1/T_f^2$缩放，而同一几何位置的重力项不随时间变化。主轨迹随执行时间增加，峰值逐渐接近约$8.135\units{N\,m}$的离散静态极限；近奇异轨迹在$16\units s$时仍需要约$347.25\units{N\,m}$，接近约$346.90\units{N\,m}$的离散静态极限。

\begin{table}[htbp]
\centering\caption{不同描绘时间下两电机的最大绝对力矩，401点扫参}\label{tab:speed}
\begin{tabular}{ccc}\toprule 执行时间/s & 主轨迹/(N$\cdot$m) & 近奇异轨迹/(N$\cdot$m)\\\midrule
2 & 8.586 & 369.500\\
4 & 8.247 & 352.551\\
8 & 8.163 & 348.314\\
16 & 8.142 & 347.254\\
静态极限 & 8.135 & 346.901\\\bottomrule\end{tabular}
\end{table}

因此，降低速度能够减小动态附加负载，却不能改变给定机构位形的静态传力关系。对需要保持笔尖接触并承受工具自重的竖直描绘任务，应优先把轮廓布置在远离并联奇异性的区域，再用速度规划调整动态余量。若仅根据“慢速运动惯性小”就允许轨迹靠近奇异位形，会遗漏本例中占主要量级的重力传力问题。
